\par\Needspace{4\baselineskip}
\begin{center}
\SourceLeaf{027b}\pdfbookmark[2]{9.}{article-009}9.
\end{center}
If the perpendicular drawn from any point of the ellipse to the line of apsides is
extended backwards until it meets the circle described about the centre of the ellipse
with radius $a$, the inclination to the line of apsides of the radius corresponding to
the point of intersection (understood in the same way as above for the true anomaly)
will equal the eccentric anomaly, as is readily deduced from equation IX of the
preceding article. It is also clear that $r\sin v$ is the distance of any point of
the ellipse from the line of apsides. Since by equation VIII this distance becomes
$= a\cos\varphi\sin E$, it is greatest for $E = 90^0$, that is, at the centre of the
ellipse. This maximum distance, which is
$= a\cos\varphi = \frac{p}{\cos\varphi} = \BareRoot{ap}$, is called the
\emph{minor semiaxis}. At the focus of the ellipse, that is, for $v = 90^0$, this
distance clearly becomes $= p$, or equal to the semiparameter.
